% Part: first-order-logic
% Chapter: tableaux
% Section: rules-and-proofs

\documentclass[../../../include/open-logic-section]{subfiles}

\begin{document}

\iftag{FOL}
      {\olfileid{fol}{tab}{rul}}
      {\olfileid{pl}{tab}{rul}}

\olsection{नियम आणि \usetoken{P}{tableau}}

!!^a{tableau} म्हणजे !!a{sentence} !!a{structure} मध्ये ज्या संभाव्य प्रकारे
सत्य अथवा असत्य असू शकते त्यांचा पद्धतशीर आढावा होय. टॅब्लोचे मूलभूत
घटक !!{signed formula}s असतात: !!{sentence}s सोबत सत्यतामूल्याचे
``चिन्ह,'' म्हणजे $\True$ किंवा~$\False$. ही चिन्हांकित !!{formula}s
(खाली वाढणाऱ्या) वृक्षात मांडली जातात.

\begin{defn}
  एक \emph{!!{signed formula}} म्हणजे सत्यतामूल्य आणि !!a{sentence} यांची जोडी,
  म्हणजे पुढीलपैकी एक:
  \[
  \sFmla{\True}{!A} \text{ किंवा } \sFmla{\False}{!A}.
  \]
\end{defn}

सहज अर्थानुसार, $\sFmla{\True}{!A}$ चे वाचन आपण ``$!A$ सत्य असू
शकते'' आणि $\sFmla{\False}{!A}$ चे वाचन ``$!A$ असत्य असू शकते''
(एखाद्या !!{structure} मध्ये) असे करू शकतो.

वृक्षातील प्रत्येक !!{signed formula} एकतर \emph{गृहीतक} असते (ती
वृक्षाच्या अगदी वरच्या टोकाला नोंदवलेली असतात), किंवा तिच्या वर असलेल्या
!!a{signed formula} ला अनेक निगमन नियमांपैकी एखादा नियम लावून ती मिळते.
पूर्ववर्ती !!{formula} च्या प्रत्येक संभाव्य !!{main operator} साठी दोन नियम
असतात: चिन्ह~$\True$ असेल त्या स्थितीसाठी एक आणि चिन्ह~$\False$ असेल त्या
स्थितीसाठी एक. काही नियम वृक्षाची शाखा फोडू देतात, तर काही नियम शाखेत फक्त
!!{signed formula}s जोडतात. एखादा नियम फक्त अगदी आधी येणाऱ्या
!!{signed formula} वरच लागू करावा असे नाही; मुळापासून नियम लागू करण्याच्या
स्थानापर्यंतच्या शाखेतील कोणत्याही !!{signed formula} वर तो लागू करता येतो
(आणि अनेकदा आधीच्या एखाद्यावर लावावाच लागतो).

एखाद्या शाखेत $\sFmla{\True}{!A}$ आणि $\sFmla{\False}{!A}$ ही दोन्ही
असतील तेव्हा ती शाखा \emph{बंद} असते. ज्याची प्रत्येक शाखा बंद असते तो
!!{tableau} बंद असतो. सहज अर्थानुसार, कोणतीही शाखा एक संयुक्त शक्यता
वर्णन करते; परंतु $\sFmla{\True}{!A}$ आणि $\sFmla{\False}{!A}$ एकाच वेळी
शक्य नाहीत. दुसऱ्या शब्दांत, एखादी शाखा बंद असेल तर तिने वर्णन केलेली शक्यता
बाद होते. विशेषतः याचा अर्थ असा की, बंद !!{tableau} हा
$\sFmla{\True}{!A}$ स्वरूपाचे प्रत्येक गृहीतक सत्य आणि
$\sFmla{\False}{!A}$ स्वरूपाचे प्रत्येक गृहीतक असत्य करणाऱ्या सर्व शक्यता
बाद करतो.

बंद !!{tableau} \emph{$!A$ साठीचा} म्हणजे ज्याचे
मूळ~$\sFmla{\False}{!A}$ आहे असा बंद !!{tableau}. असा बंद !!{tableau}
अस्तित्वात असल्यास, $!A$ असत्य असण्याच्या सर्व शक्यता बाद होतात; म्हणजेच
$!A$ प्रत्येक !!{structure} मध्ये सत्य असले पाहिजे.

\end{document}
